\section{Introduction}
\label{sec:introduction}

Problem~D concerns the screening of inorganic perovskite compositions
against five material requirements: a formation-energy limit, a target
band-gap interval, and two band-edge constraints.  Evaluating these
properties at scale can be costly, while the candidate space expands with
the choices of A-site element, B-site element, and anion group.  The task
therefore combines two steps: predict a composition-level violation score,
then select candidates under fixed- and variable-anion settings
\cite{apmcm2026problemD}.

Composition-based features and supervised learning offer a practical
first-pass screen when structural calculations are unavailable
\cite{Ward2016Magpie,Chen2024PerovskiteML}.  Average prediction error alone,
however, does not determine whether the model is useful for this task.  With
a limited verification budget, the model must place low-violation
compositions near the top of its ranking.  Global accuracy and low-tail
retrieval can conflict, and a high overall rank correlation does not
guarantee reliable ordering among the best-scoring candidates.  We therefore
use the surrogate to prioritize follow-up calculations, not to certify
physical feasibility.

QUBO and Ising formulations express discrete selection problems in a
form suitable for annealing-based solvers
\cite{Lucas2014IsingNP,Glover2019QUBOTutorial}.  The conversion introduces
a separate numerical question: a floating-point objective may lose
feasibility or decision-relevant score differences when mapped to a finite
coefficient range.  Large constraint penalties can protect one-hot
feasibility but leave less resolution for the material objective
\cite{Ceselli2023GoodQUBOEncodings,Muecke2025OptimumPreservingQUBO}.
Solver samples are therefore interpreted only after checking the encoding,
constant offset, coefficient range, one-hot constraints, and decoded
original-scale score.

\subsection{Research framework}

The three questions form a linked workflow.  Q1 reconstructs the five
official violation components from predicted material properties.  Its
two-stage band-gap model combines the probability of a positive gap with
the predicted gap conditional on positivity; model selection considers
both global error and low-tail retrieval.  For Q2, we fix
$X=\mathrm{NO}_2$ and represent the 2,704 ordered cation pairs with 52
A-site and 52 B-site one-hot variables.  For Q3, we solve the same
104-variable formulation separately for $\mathrm{O}_3$, $\mathrm{NO}_2$,
and $\mathrm{FO}_2$, then compare the three decoded winners using the
original surrogate scores.  The three groupwise searches cover all 8,112
Q3 candidates.  This design avoids directly quantizing the 208-variable
joint control, whose precision problems are examined below.

We check each stage against records appropriate to that stage.  Source
hashes and split assignments document the data pipeline; exhaustive
enumeration checks the legal QUBO states; and the QUBO--Ising conversion is
verified with its constant offset restored.  Coefficient range, decoded
feasibility, original-score regret, and solver sampling are assessed
separately.  The formal local solves use Kaiwu SDK~1.4.3 on a CPU, whereas
four final 104-variable matrices were submitted separately to SPQC-550.
Three platform tasks returned at least one group proxy argmin, while the
Q3/NO2 task returned no legal sample.  These small, unequal hardware
returns are not pooled with the CPU results or with the historical
17- and 30-variable short-list tasks.

\subsection{Contributions}

The evidence-supported contributions are:

\begin{enumerate}[leftmargin=2em]
 \item an auditable composition-only surrogate that explicitly balances
       overall error and low-tail screening, with grouped validation,
       partition-sensitivity checks, and an exact finite-budget random baseline;
 \item exhaustive proxy-score evaluation for Q2 and groupwise Q3, followed
       by one-hot encoding and decoding on the original score scale;
 \item an oracle-assisted, group-winner-preserving two-level integer
       encoding that fits the recorded eight-bit coefficient range, while
       explicitly sacrificing non-winner score differences and complete
       ranking information;
 \item documented comparisons of CPU-local SDK and simulated-annealing
       samples, including feasibility, optimum hits, timing records,
       negative tests, and regression checks; and
 \item a four-layer error analysis separating surrogate, mathematical
       mapping, coefficient precision, and solver sampling uncertainty.
\end{enumerate}

Together, these components form a screening and optimization workflow
for the supplied candidate domain.  The study does not establish the
physical feasibility of the selected materials, repeatable hardware
performance, or quantum advantage.